\documentclass[english,utf8,aspectratio=169]{beamer}
\usetheme{tud}
\usecolortheme{tud}
\usepackage[backend = biber,style = numeric]{biblatex}
\addbibresource{references.bib}
\usepackage[inkscapepath={./build/svg-inkscape}]{svg}
\svgpath{{./img/}}

\AtBeginSection[]{\partpage{\usebeamertemplate***{part page}}}
\setbeamertemplate{theorems}[numbered]

\title{Lock-Free Deduplication}
\subtitle{Implementation considerations}
\author{Eric Niklas Wolf}
\institut{Chair of Distributed and Networked Systems}

\begin{document}
    \maketitle

    \begin{frame}
        \frametitle{Outline}
        \tableofcontents
    \end{frame}

    \section{Introduction}
    \begin{frame}
        \frametitle{Fossil collection step}
        \begin{columns}
            \begin{column}{.5\textwidth}
                \begin{itemize}
                    \item identifies chunks to be deleted
                    \item archives may be created in parallel, possibly after
                        downloading manifest files
                \end{itemize}
                To prevent a race condition causing chunks to be deleted prematurely they
                are first turned into fossils.
            \end{column}
            \begin{column}{.5\textwidth}
                \includesvg[height=0.8\textheight, width=0.9\textwidth]{fossil_collection.drawio}
            \end{column}
        \end{columns}
    \end{frame}

    \begin{frame}
        \frametitle{Fossil deletion step}
        \begin{columns}
            \begin{column}{.5\textwidth}
                \begin{itemize}
                    \item uses an existing fossil collection
                    \item fossils which where referenced in the meantime are recovered
                \end{itemize}
                Is delayed~\cite[Policy 3]{Duplicacy} until referencing a chunk would cause a new
                one to be uploaded~\cite[Policy 2]{Duplicacy}.
            \end{column}
            \begin{column}{.5\textwidth}
                \includesvg[height=0.8\textheight, width=0.9\textwidth]{fossil_deletion.drawio}
            \end{column}
        \end{columns}
    \end{frame}

    \section{Non-atomic retrieval operation}
    \begin{frame}
        \frametitle{Accessing all archives}
        \begin{itemize}
            \item Lock-Free Deduplication requires an operation to retrieve all manifest
                files~\cite[Operation `downloadAllManifestFiles']{Duplicacy}
            \item on backends like POSIX filesystems this operation is split into listing
                all files and accessing them
            \item changes to a directory during listing may not be picked up
            \item accessing files may itself be a non-atomic operation (chunked manifests)
            \item since manifest files can be deleted at any time by a concurrently running
                pruning operation accessing listed files can fail because they where deleted
        \end{itemize}
    \end{frame}

    \begin{frame}
        \frametitle{Solution}
        \begin{lemma}
            Whether manifest files created or deleted during retrieval are picked up by the operation
            or not does not disrupt the Lock-Free Deduplication algorithm.
        \end{lemma}
        \begin{corollary}
            Since ignoring deleted files after listing is equivalent to their deletion being picked
            up by the retrieval operation it is safe to do so.
        \end{corollary}
    \end{frame}

    \begin{frame}
        \frametitle{Proof idea}
        \begin{columns}
            \begin{column}{.6\textwidth}
                The result of an atomic retrieval operation is a snapshot of a point during its
                execution while the non-atomic one allows creations and deletions happening during
                its execution to be moved freely before or after this point.
            \end{column}
            \begin{column}{.4\textwidth}
                \includesvg[height=0.8\textheight, width=0.9\textwidth, pretex=\small]{retrieval_cases.drawio}
            \end{column}
        \end{columns}
        \vfill
        It will be shown that any guarantees of an atomic implementation used by the algorithm
        are also provided by an non-atomic operation.
    \end{frame}

    \begin{frame}
        \frametitle{Proof}
        \begin{proof}
            If a manifest operation happens during the execution of the retrieval operation an atomic
            implementation can't guarantee whether or not the operation happens before the snapshot
            or after it, the same as a non-atomic operation.
            The same is true for manifest operations happening before or after the execution of the
            retrieval operation, since a non-atomic operation can not move them before or after the
            snapshot giving them the same visibility when using an atomic operation.
            The only difference is that the non-atomic implementation does not preserves the ordering
            between different manifest operations happening during the execution of the retrieval operation,
            which is not used by the algorithm and it is assumed the clients won't either.
        \end{proof}
    \end{frame}

    \begin{frame}
        \frametitle{On manifest operation ordering}
        The fossil deletion step assumes that after Policy 3 is satisfied the retrieval operation
        will have seen any archive which may reference its fossils, leading to problems in the
        following scenario should only the second manifest be seen:
        \begin{center}
            \includesvg[height=0.6\textheight, width=0.9\textwidth, pretex=\small]{manifest_ordering.drawio}
        \end{center}
    \end{frame}

    \begin{frame}
        \frametitle{Solution}
        \begin{theorem}
            Performing a second listing operation after Policy 3 is satisfied will see any manifests
            which may reference the fossils.
        \end{theorem}
        \begin{proof}
            Assuming an manifest which references the fossils is not seen by the retrieval operation.
            Since it has to be created before the manifest satisfying Policy 3 for its
            client~\cite[Policy 2]{Duplicacy} it has to be created during the retrieval operation
            and therefore before a following second listing operation.
        \end{proof}
    \end{frame}

    \section{Non-atomic chunk access}
    \begin{frame}
        \frametitle{Accessing a chunk}
        \begin{block}{Fossil access policy~\cite[Policy 2]{Duplicacy}}
            \begin{enumerate}
                \item A restore, list, or check procedure that reads existing backups
                    can read the fossil if the original chunk cannot be found.
                \item A backup procedure cannot access any fossils.
                    That is, it must upload a chunk if it cannot find the chunk,
                    even if a corresponding fossil exists.
            \end{enumerate}
        \end{block}
        \begin{theorem}
            The fossil access policy is affected by a race condition in case the chunk and fossil
            accesses are two separate operations which can cause chunks to appear as missing.
        \end{theorem}
    \end{frame}

    \begin{frame}
        \frametitle{Possible fossil operations}
        Archive \#1 and \#2 reference the same chunk (exclusively).
        \begin{center}
            \includesvg[height=0.75\textheight, width=0.9\textwidth, pretex=\small]{nonatomic_access.drawio}
        \end{center}
    \end{frame}

    \begin{frame}
        \frametitle{Proof}
        \begin{proof}
            When executing the fossil collection and deletion steps started in the example the
            following sequence of events is possible:
            \begin{enumerate}
                \item A chunk is turned into a fossil (by a fossil collection step)
                \item A client tries to access the chunk and fails
                \item The chunk is recovered from its fossil (by a fossil deletion step)
                \item The client tries to access the fossil and fails
            \end{enumerate}
            From the perspective of the client the chunk is missing from the repository.
            The sequence can be repeated as long as there are pending fossil deletions.
        \end{proof}
    \end{frame}

    \begin{frame}
        \frametitle{Solution}
        A client may alternate between trying to access a chunk and its fossil, but can this
        action result in an infinite loop?
        \begin{theorem}
            Assuming every chunk access is tied to an archive which references the chunk, the
            accessing client generates a finite amount of new archives while accessing it
            and clients only have a finite amount of pending fossil deletion steps which in turn
            only perform a finite amount of operations on each fossil.
            By periodically checking whether this archive exists and references the chunk
            a client will after a finite amount of time either find the chunk or detect that it
            was deleted.
        \end{theorem}
    \end{frame}

    \begin{frame}
        \frametitle{Proof}
        \begin{lemma}
            \label{lemma:finiteoperations}
            Under the stated assumptions a client can assume that a fossil it is trying to access
            will only be operated on by a finite number of operations.
        \end{lemma}
        \begin{proof}
            When the client creates only a finite amount there has to be an archive which is the
            last archive created during the chunk access by the client.
            A fossil deletion step must make sure Policy 3~\cite{Duplicacy} is satisfied before
            continuing, causing new pending fossil deletion steps to wait until the client creates
            a new archive.
            The existing finite number of pending fossil deletion steps can together only perform
            a finite amount of operations on the fossil.
        \end{proof}
    \end{frame}

    \begin{frame}
        \frametitle{Proof}
        \begin{proof}
            After the finite amount of operations in Lemma~\ref{lemma:finiteoperations} have passed
            the client may either find the chunk or its fossil or notice the archive check failing.
            The archive check succeeding implies the referenced chunk (or fossil) existed at the time
            of the check, which combined with no new operations on its fossil will prevent it from
            being deleted in the meantime therefore making the check succeeding impossible should
            this be the case.
            %NOTE: checking the archive existence alone is not enough since it can be replaced by
            %another archive referencing different chunks.
        \end{proof}
    \end{frame}

    \section{Race condition during fossil collection and deletion}
    \begin{frame}
        \frametitle{Race condition during fossil collection and deletion}
        \begin{theorem}
            It is possible to delete a chunk while it is still being referenced by an archive
            when following the Lock-Free Deduplication algorithm.
        \end{theorem}
    \end{frame}

    \begin{frame}
        \frametitle{Proof}
        \begin{proof}
            Assuming a repository which has 3 clients, one of them being about to perform a
            fossil deletion after all of them created archives to satisfy Policy 3.
            After this fossil deletion downloads all archive files to check Policy 3 and determining
            a certain fossil to be safe to delete the other two clients beform the following actions:
            \begin{enumerate}
                \item create an archive which causes the chunk of the fossil to be recreated
                \item initiate a fossil collection to prune this archive by downloading all archive files
                \item create another archive which also references the chunk of the first archive
                \item finish the fossil collection by turning the recreated chunk into a fossil
            \end{enumerate}
        \end{proof}
    \end{frame}

    \begin{frame}
        \frametitle{Proof}
        \begin{proof}
            Now the second archive references a fossil which would be recovered by the following
            fossil deletion step.
            But since the already running fossil deletion step determined the same fossil to be safe
            to delete it will delete it causing the second archive to reference a not existing chunk.
        \end{proof}
    \end{frame}

    \begin{frame}
        \frametitle{Solution}
        \begin{theorem}
            \label{theorem:uniquefossils}
            By preventing fossils produced by a fossil collection from being deleted by a fossil
            deletion not associated with it referenced chunks can't be deleted.
        \end{theorem}
    \end{frame}

    \begin{frame}
        \frametitle{Proof}
        \begin{proof}
            Assuming a referenced chunk was deleted by a fossil deletion step.
            This means the chunk was turned into a fossil by the associated fossil collection step,
            implying all archives referencing it where selected to be pruned.
            But it being referenced after the fossil deletion step means the creation of a new archive
            not seen by the fossil collection step referenced the chunk before Policy 2 would have
            caused it to be recreated, otherwise the recreated chunk would have caused its fossil
            to be safe to delete.
            Since such an archive would have caused the check for Policy 3 to fail until it was seen
            by the fossil deletion step it would have caused the chunk to not be deleted.
        \end{proof}
    \end{frame}

    \section{Recovering from a failed archive creation}
    \begin{frame}
        \frametitle{Archive creation failure}
        When creating an archive fails after some chunks where already uploaded to the repository
        they are not deleted until another archive which references them is created and pruned,
        which may never happen.
        \begin{theorem}
            \label{Theorem:archivefailure}
            It is possible to perform a fossil collection step which will identify these chunks
            as not being referenced while not disrupting the Lock-Free Deduplication algorithm,
            assuming there is a way to enumerate all chunks existing within a repository.
        \end{theorem}
    \end{frame}

    \begin{frame}
        \frametitle{Proof}
        \begin{lemma}
            \label{lemma:archiveremoval}
            Removing an archive does not disrupt the Lock-Free Deduplication algorithm besides
            leaving unreferenced chunks behind.
        \end{lemma}
        \begin{proof}
            An archive being removed before being seen by a fossil collection will at most cause
            it to not turn some chunks into fossils, therefore not deleting them during fossil deletion.
            An archive being removed after being seen by a fossil collection will at most cause its
            fossil deletion to be delayed until the client creates a new archive.
        \end{proof}
    \end{frame}

    \begin{frame}
        \frametitle{Proof}
        \begin{proof}
            Assuming an archive referencing all chunks existing in the repository at a given point
            was created, pruning it would take care of all chunks which are not referenced by any
            other archive.
            Since we can safely pretend this archive was immediately removed after being created as
            proven by Lemma~\ref{lemma:archiveremoval} and therefore not seen by any other client we
            can skip creating such an archive and just perform the fossil collection associated with
            its pruning operation, treating the set of all chunks existing within the repository as
            the set of chunks referenced by it.
            Since this imaginary archive did not create any new chunks there are no new unreferenced
            chunks to worry about.
        \end{proof}
    \end{frame}

    \section{Recovering from a failed fossil collection}
    \begin{frame}
        \frametitle{Fossil collection failure}
        When performing a fossil collection fails after some chunks where already turned into fossils
        they must be prevented from existing in the repository indefinitely.
        \begin{theorem}
            Such fossils can be safely removed by first recovering all fossils existing within a
            repository followed by the procedure described in Theorem~\ref{Theorem:archivefailure}.
        \end{theorem}
    \end{frame}

    \begin{frame}
        \frametitle{Proof}
        \begin{proof}
            Idea: no chunks are deleted, erroneously recovered ones are taken care of by
            Theorem~\ref{Theorem:archivefailure}.
        \end{proof}
    \end{frame}

    \section{Adding clients}
    \begin{frame}
        \frametitle{Adding clients}
        The paper~\cite{Duplicacy} treats the set of clients using a repository as being static,
        but can this set be expanded?
        \begin{theorem}
            Expanding the set of clients does not disrupt the Lock-Free Deduplication algorithm. 
        \end{theorem}
    \end{frame}

    \begin{frame}
        \frametitle{Proof}
        \begin{lemma}
            \label{lemma:runafteradd}
            Adding a client will not disrupt future fossil collections and deletions.
        \end{lemma}
        \begin{proof}
            Adding a client will cause future fossil collections and deletions to treat him the same
            as the existing clients, behaving as if the client has been part of a static client set
            since the repository was created.
        \end{proof}
    \end{frame}

    \begin{frame}
        \frametitle{Proof}
        \begin{lemma}
            \label{lemma:collectwhileadd}
            Adding a client will not disrupt fossil collections which where running at the time of
            his addition.
        \end{lemma}
        \begin{proof}
            The only way a new client can interact with a running fossil collection step is by
            creating a new archive.
            Should the fossil collection step see the new archive it will treat it the same as archives
            created by existing clients, behaving as if the client has been part of a static client
            set since the repository was created.
            Should the fossil collection step not see the new archive it will behave the same as if
            the client had been part of a static client set since the repository was created, with
            the client not having created a single archive yet.
        \end{proof}
    \end{frame}

    \begin{frame}
        \frametitle{Proof}
        \begin{lemma}
            \label{lemma:deletewhileadd}
            Adding a client will not disrupt fossil deletions which where running at the time of his
            addition.
        \end{lemma}
        \begin{proof}
            The fossil deletion seeing the new client will cause it to behave as if the client had
            been part of a static client set since the repository was created, otherwise it will
            behave as if the client does not exist.
            Since a fossil deletion implementing Theorem~\ref{theorem:uniquefossils} will only access
            fossils created by its associated fossil collection it won't interact the chunks created
            by the new client since they have to be recreated should the client attempt to create archives
            referencing these fossils according to Policy 2.
        \end{proof}
    \end{frame}

    \begin{frame}
        \frametitle{Proof}
        \begin{lemma}
            \label{lemma:usebeforeadd}
            Adding a client will not disrupt fossil deletions which where using fossils created by
            fossil collections which finished before it.
        \end{lemma}
        \begin{proof}
            In case a fossil deletion is running at the time a client is being added it will not
            be disrupted as proven by Lemma~\ref{lemma:deletewhileadd}.
            Should this not be the case it will behave the same as if the client had been part of a
            static client set since the repository was created, with the client not having created
            a single archive yet.
        \end{proof}
    \end{frame}

    \begin{frame}
        \frametitle{Proof}
        \begin{proof}
            The possible interactions between the addition of a new client and the Lock-Free
            Deduplication algorithm are as follows:
            \begin{enumerate}
                \item a fossil collection or deletion was running after it, which is safe as proven
                    by Lemma~\ref{lemma:runafteradd}
                \item a fossil collection or deletion step was running during it, which is safe as proven
                    by Lemma~\ref{lemma:collectwhileadd} and Lemma~\ref{lemma:deletewhileadd}
                \item a fossil deletion step is running after its associated fossil collection step
                    finished before it, which is safe as proven by Lemma~\ref{lemma:usebeforeadd}
            \end{enumerate}
            In all cases adding a client won't disrupt the Lock-Free Deduplication algorithm.
        \end{proof}
    \end{frame}

    \section{Removing clients}
    \begin{frame}
        \frametitle{Removing clients}
        The paper~\cite{Duplicacy} treats the set of clients using a repository as being static,
        but can this set be reduced?
        \begin{theorem}
            Reducing the set of clients does not disrupt the Lock-Free Deduplication algorithm beyond
            the creation of unreferenced chunks and fossils without an associated fossil collection.
        \end{theorem}
    \end{frame}

    \begin{frame}
        \frametitle{Proof}
        \begin{proof}
            Idea: The removal of clients only affects the checking for Policy 3, which is not a
            problem since the policy ignoring a client can only happen after it finished any archive
            creation which would result in the same outcome.
        \end{proof}
    \end{frame}

    \section{References}
    \begin{frame}
        \frametitle{References}
        \printbibliography{}
    \end{frame}
\end{document}
